\documentclass[10pt,a4paper]{amsart}
\usepackage[margin=19mm]{geometry}
\usepackage{amsmath,amssymb,mathtools}
\usepackage[scr=rsfs,cal=euler]{mathalfa}
\usepackage{xfrac}
\usepackage[unicode,hidelinks,bookmarks=false]{hyperref}
\newcommand{\ie}{i.e.\ }
\newcommand{\cf}{cf.\ }
\newcommand{\resp}{respectively\ }
\newcommand{\Subsection}{\subsection}
\providecommand{\nobreakdash}{\nobreak\hbox{-}}
\setlength{\emergencystretch}{2em}
\multlinegap0pt
\usepackage{bm}
\usepackage{bbm}
\usepackage{dsfont}
\usepackage{xspace}
\usepackage{cancel}
\usepackage{mathtools}
\usepackage[shortlabels]{enumitem}
\usepackage{amsthm}
\makeatletter
\@ifundefined{theorem}{\newtheorem{theorem}{Theorem}}{}
\@ifundefined{corollary}{\newtheorem{corollary}[theorem]{Corollary}}{}
\@ifundefined{remark}{\newtheorem{remark}[theorem]{Remark}}{}
\makeatother
\def\bigmid{\ \rule[-3.5mm]{0.1mm}{9mm}\ }
\newcommand{\ZZ}{{\mathbb Z}}
\title[Barcode entropy and relative symplectic cohomology]{Barcode entropy and relative symplectic cohomology}
\author{Jonghyeon Ahn}
\date{Source: arXiv:2601.15606v1 (CC BY 4.0)}
\begin{document}
\maketitle

\noindent\emph{Source and licence.} Barcode entropy and relative symplectic cohomology. Version arXiv:2601.15606v1 (CC BY 4.0), distributed under the Creative Commons Attribution 4.0 International licence, as recorded in the arXiv licence metadata for this exact version and shown on the abstract page https://arxiv.org/abs/2601.15606v1. Full rights notice: licenses/arxiv-2601.15606-CC-BY-4.0.txt.\par\smallskip

\noindent\textit{Editorial scope.} Counted unit: the Theorem whose source label is \texttt{rene} in the article (source0.tex lines 1224--1235, source label \texttt{rene}), delivered together with the article's own complete original proof of it (source0.tex lines 1236--1288, one \texttt{proof} environment of 53 lines). No mathematics is written, repaired or re-derived: statement and proof are the article's own text line for line. Required context retained uncounted: chsh (lines 600--602); ua (lines 723--725); bound (lines 782--784); hamiso (lines 986--992); hamisocoro (lines 1099--1105); rmk1 (lines 1110--1137). Every definition, display and label of those lines is retained unchanged. Omitted from this package and not redistributed: 1--599, 603--722, 726--781, 785--985, 993--1098, 1106--1109, 1138--1223, 1289--1485. Nothing inside a retained excerpt is omitted or altered: every retained body is delivered exactly as the article's lines are written, so no omission receipt of any kind accompanies this package. Citations: the retained excerpts carry no citation command at all, so no bibliography entry of the article is transcribed and no reference is redistributed. Reference adaptation: none was needed and none was made; every reference inside the delivered unit resolves inside it, so no unresolved reference or citation is produced. Printed slips kept literal: no printed word, symbol, hypothesis or number of the counted statement or of its proof was altered. Layout and class: the article's own class is \texttt{amsart} with packages including graphicx, amssymb, latexsym, eucal, hyperref, geometry, inputenc, amsmath, amsthm, bigints, amsfonts, epsfig, nicefrac, xy; the wrapper is the lane's pinned \texttt{amsart} editorial wrapper at 10pt on A4, which restates the theorem-like environments the retained text needs and adds the article's own macro definitions that the retained text needs (\texttt{\textbackslash ZZ}, \texttt{\textbackslash bigmid}), with extra wrapper packages (bm, bbm, dsfont, xspace, cancel, mathtools, enumitem). The delivered numbering is the unit's own. Assets: no retained span contains a figure environment, a graphics-inclusion command, a PDF-page inclusion, a table or a TikZ picture; no image file is copied into this package, the package declares no asset dependency and no image file is redistributed with it. Licence: version arXiv:2601.15606v1 of this article is distributed under the Creative Commons Attribution 4.0 International licence, as recorded in the arXiv licence metadata for this exact version and shown on the abstract page https://arxiv.org/abs/2601.15606v1. A full rights notice is delivered at licenses/arxiv-2601.15606-CC-BY-4.0.txt. Characters: every retained excerpt is pure ASCII, so the delivered engine, XeLaTeX with its default Latin Modern text font, reports no missing character.
\begin{align}\label{chsh}
    c_H \leq 3r_1 s_H.
\end{align}

\begin{align}\label{ua}
    \mathcal{A}_H(x) \geq c_H-(1+3r_1) \mathcal{A}(\gamma).
\end{align}

    \begin{align}\label{bound}
       -1-r_1 \leq \frac{d}{d\tau}a_{h_1}(\tau) \leq -1.
    \end{align}

\begin{theorem}\label{hamiso}
For each integer $k \in \ZZ$, there exists a contact type $K$-semi-admissible Hamiltonian function $H$ such that 
    \begin{align*}
        SH^{>-\tau,k}_M(K) \cong HF^{(a_{H}(\tau),0],k}(H)
    \end{align*}
    for $0 \leq \tau \leq s_H$.
\end{theorem}

\begin{corollary}\label{hamisocoro}
    Let $H$ be a $K$-semi-admissible Hamiltonian function. Then for each $k \in \ZZ$, there exists a constant $\sigma >0$ such that
    \begin{align*}
        SH^{>-\tau,k}_M(K) \cong HF^{(a_{\sigma H}(\tau),0],k}(\sigma H)
    \end{align*}
    for $0 \leq \tau \leq s_{\sigma H} = \sigma s_H$.
\end{corollary}

\begin{remark}\label{rmk1}
The proofs of Theorem \ref{hamiso} and Corollary \ref{hamisocoro} also imply the following statements.
\begin{enumerate}[label=(\alph*)]
    \item Let $A \subset \ZZ$ be a finite subset of $\ZZ$. \vspace{0.2cm}
    \begin{itemize}
       \item There exists a contact type $K$-semi-admissible Hamiltonian function $H$ such that there exists an isomorphism
    \begin{align*}
        SH^{>-\tau,k}_M(K) \cong HF^{(a_{H}(\tau),0],k}(H)
    \end{align*}
    for $0 \leq \tau \leq s_H$ and for $k \in A$. \vspace{0.2cm}
    \item For any contact type $K$-semi-admissible Hamiltonian function $H$, there exists constant $\sigma >0$ such that
    \begin{align*}
         SH^{>-\tau,k}_M(K) \cong HF^{(a_{\sigma H}(\tau),0],k}(\sigma H)
    \end{align*}
    for $0 \leq \tau \leq s_{\sigma H}$ and for $k \in A$.\vspace{0.2cm}
    \end{itemize}
    \item Let $H$ be a contact type $K$-semi-admissible Hamiltonian function and let
    \begin{align*}
        I_H = \left\{ k \in \ZZ \bigmid \substack{ \textstyle \text{the actions of upper orbits of } H \text{ of Conley-Zehnder indices } $\vspace{0.2cm}$\\ 
             \,\,\textstyle k-1, k \text{ and } k+1 \text{ are positive}} \right\}.
    \end{align*}
    Then there exists an isomorphism 
    \begin{align*}
        SH_M^{>-\tau,k}(K) \cong HF^{(a_{H}(\tau),0],k}(H).
    \end{align*}
    for $0\leq\tau\leq s_H$ and for $k \in I_H$.
   \end{enumerate}
       \end{remark}

\begin{theorem}\label{rene}
    There exists a contact type $K$-semi-admissible Hamiltonian function $H$ such that for any $\epsilon>0$, 
         \begin{align*}
       \hbar_{\epsilon}(SH_M(K)) \leq \hbar_{\epsilon}(HF^{\textrm{neg}}(H))
    \end{align*}
    and hence, taking $\epsilon \to 0$, we have  
    \begin{align*}
       \hbar(SH_M(K)) \leq \hbar(HF^{\textrm{neg}}(H)).
        \end{align*}
  
    
\end{theorem}

\begin{proof}
Choose a $K$-semi-admissible Hamiltonian function $H$ satisfying
\begin{align}\label{ch}
    c_H \geq 1+3r_1
\end{align}
Observe that the condition \eqref{ch} only depends on $r_1$, that is, on the position of $K$ inside $M$ and it forces to $s_H$ to be greater than 1. Indeed, if $s_H \leq 1$, then
\begin{align*}
    1+3r_1 &\leq c_H &&\textrm{assumption}\\
    &\leq 3r_1 s_H &&\textrm{by \eqref{chsh}}\\
    &\leq 3r_1,
\end{align*}
which yields a contradiction.

\vspace{0.2cm}

For each $\sigma>0$, the number of bars $b_{\epsilon} \left( \textrm{tru}(SH_M(K),\sigma ) \right)$ is finite and hence only finitely many grading degrees contribute nontrivially to it. Consequently, there exists a finite subset $A_\sigma$ of $\ZZ$ such that
\begin{align*}
    b_{\epsilon} \left( \textrm{tru}(SH_M(K),\sigma ) \right) = b_{\epsilon} \left( \textrm{tru}\left(\bigoplus_{k \in A_\sigma}SH^k_M(K),\sigma \right) \right) = \sum_{k \in A_\sigma}b_{\epsilon} \left( \textrm{tru}(SH^k_M(K),\sigma ) \right).
\end{align*}
To estimate $b_{\epsilon} \left( \textrm{tru}(SH_M(K),\sigma ) \right)$, it suffices to consider Reeb orbits of $(\partial K,\alpha)$ with periods less than $\sigma - \epsilon$. We claim that
\begin{align}\label{ingiso}
     SH^{>-\tau,k}_M(K) \cong  HF^{(a_{\sigma H}(\tau),0],k}(\sigma H)
\end{align}
for $0\leq \tau\leq\sigma \leq \sigma s_H$ and $k \in A_\sigma$. To justify this claim, it is enough to show that the upper orbits of $\sigma H$ have positive actions in view of the proof of Theorem \ref{hamiso}, Corollary \ref{hamisocoro} and Remark \ref{rmk1}. Let $x$ be an upper orbit of $\sigma H$ corresponding to a Reeb orbit $\gamma$ with $\mathcal{A}(\gamma) < \sigma - \epsilon$. Then we obtain
\begin{align*}
   \mathcal{A}_{\sigma H}(x) &\geq c_{\sigma H} - (1+3r_1)\mathcal{A}(\gamma) &&\textrm{by \eqref{ua}}\\&> \sigma c_H- (1+3r_1)(\sigma -\epsilon)\\&>0 &&\textrm{by the choice of }\,\,H,
\end{align*}
and this proves the claim.

\vspace{0.2cm}

Recall from \eqref{bound} that for $0 \leq \tau_1 \leq \tau_2 \leq \sigma \leq \sigma s_H = s_{\sigma H}$,
\begin{align*}
\tau_2 - \tau_1 \leq a_{\sigma H}(\tau_1) - a_{\sigma H}(\tau_2) \leq (1+r_1)(\tau_2 - \tau_1),
\end{align*}
and this together with the isomorphism \eqref{ingiso} implies that 
\begin{align}\label{oy}
    b_{\epsilon} \left( \textrm{tru}\left( \bigoplus_{k \in A_\sigma} SH^{k}_M(K),\sigma \right) \right) \leq b_{\epsilon} \left( \textrm{tru}\left(\bigoplus_{k \in A_\sigma}HF^{\textrm{neg},k}(\sigma H),-a_{\sigma H}(\sigma)\right) \right).
\end{align}
Therefore,
\begin{align*}
      \hbar_\epsilon(SH_M(K)) &= \limsup_{\sigma \to \infty}\frac{1}{\sigma} \log^+b_{\epsilon} \left( \textrm{tru}(SH_M(K),\sigma) \right)\\ 
      &=\limsup_{\sigma \to \infty}\frac{1}{\sigma} \log^+b_{\epsilon} \left( \textrm{tru}\left(\bigoplus_{k \in A_\sigma}SH^k_M(K),\sigma)\right) \right)
      \\& \leq\limsup_{\sigma \to \infty}\frac{1}{\sigma} \log^+b_{\epsilon}\left( \textrm{tru}\left(\bigoplus_{k \in A_\sigma}HF^{\textrm{neg},k}(\sigma H),-a_{\sigma H}(\sigma)\right) \right)&&\textrm{by}\,\,\eqref{oy}\\
      & \leq\limsup_{\sigma \to \infty}\frac{1}{\sigma} \log^+b_{\epsilon}\left( \textrm{tru}\left(\bigoplus_{k \in A_\sigma}HF^{\textrm{neg},k}(\sigma H),\sigma i_H\right) \right)&&-a_{\sigma H}(\sigma) \leq  i_{\sigma H}\\
      &\leq\limsup_{\sigma \to \infty}\frac{1}{\sigma} \log^+b_{\epsilon}\left( \textrm{tru}\left(HF^{\textrm{neg}}(\sigma H),\sigma i_H\right) \right)\\
      &= \hbar_{\epsilon}(HF^{\textrm{neg}}(H)).
\end{align*}
This completes the proof.



\end{proof}

\end{document}
